\section{Năng lượng nhiên liệu và các bài toán thực tiễn}

Phần năng lượng trong đề tuyển sinh thường không dừng ở việc tính \(Q\). Dữ kiện có thể gắn với quãng đường xe chạy, chi phí, lượng khí thải hoặc lựa chọn nhiên liệu. Cách làm ổn định nhất là xác định rõ \textbf{đại lượng so sánh} và lập chuỗi tính toán theo từng bước.

\subsection{Năng lượng theo khối lượng và theo thể tích}

Hai nhiên liệu có thể có nhiệt trị theo khối lượng khác nhau, đồng thời có khối lượng riêng khác nhau. Vì vậy, nếu đề hỏi ``cùng một lít'' thì cần đổi thể tích sang khối lượng trước.

\begin{vidu}
Nhiên liệu A có nhiệt trị \SI{44}{\mega\joule\per\kilo\gram}, khối lượng riêng \SI{0.74}{\kilo\gram\per\liter}. Nhiên liệu B có nhiệt trị \SI{46}{\mega\joule\per\kilo\gram}, khối lượng riêng \SI{0.68}{\kilo\gram\per\liter}. So sánh năng lượng tỏa ra khi đốt \SI{10}{\liter} mỗi loại.
\end{vidu}
\begin{loigiai}
Với A:
\[
m_A=10\times0.74=\SI{7.4}{\kilo\gram},
\qquad
Q_A=7.4\times44=\SI{325.6}{\mega\joule}.
\]
Với B:
\[
m_B=10\times0.68=\SI{6.8}{\kilo\gram},
\qquad
Q_B=6.8\times46=\SI{312.8}{\mega\joule}.
\]
Mặc dù B có nhiệt trị theo khối lượng lớn hơn, cùng \SI{10}{\liter} thì A cho nhiều năng lượng hơn do có khối lượng riêng lớn hơn.

\textbf{Vậy,} \SI{10}{\liter} nhiên liệu A tỏa \SI{325.6}{\mega\joule}, lớn hơn \SI{312.8}{\mega\joule} của nhiên liệu B.
\end{loigiai}

\subsection{Năng lượng và quãng đường}

Nếu xe cần \(q\) kJ năng lượng có ích cho mỗi kilomet thì:
\[
s=\frac{Q_{\text{ích}}}{q}.
\]

\begin{vidu}
Một xe dùng \SI{4.0}{\liter} nhiên liệu có \(D=\SI{0.75}{\kilo\gram\per\liter}\) và nhiệt trị \SI{44}{\mega\joule\per\kilo\gram}. Hiệu suất biến đổi năng lượng của động cơ là \(30\%\). Xe cần trung bình \SI{600}{\kilo\joule} năng lượng cơ học để đi \SI{1}{\kilo\metre}. Tính quãng đường tối đa theo mô hình trên.
\source{Phỏng theo dạng bài nhiên liệu E5/E10 trong đề tuyển sinh chuyên Hóa Hà Nội 2025 và Lâm Đồng 2026.}
\end{vidu}
\begin{loigiai}
Khối lượng nhiên liệu:
\[
m=4.0\times0.75=\SI{3.0}{\kilo\gram}.
\]
Nhiệt lượng tỏa ra:
\[
Q=3.0\times44=\SI{132}{\mega\joule}
=\SI{132000}{\kilo\joule}.
\]
Năng lượng cơ học có ích:
\[
Q_{\text{ích}}=132000\times0.30=\SI{39600}{\kilo\joule}.
\]
Quãng đường:
\[
s=\frac{39600}{600}=\SI{66}{\kilo\metre}.
\]

\textbf{Vậy,} xe đi được tối đa khoảng \SI{66}{\kilo\metre} theo các dữ kiện đã cho.
\end{loigiai}

\subsection{Năng lượng, chi phí và lựa chọn nhiên liệu}

Để so sánh kinh tế, không nên chỉ nhìn giá mỗi kilogram hoặc mỗi lít. Cần quy về cùng một lượng năng lượng có ích:
\[
\text{chi phí trên một đơn vị năng lượng có ích}
=
\frac{\text{giá tiền}}{Q_{\text{ích}}}.
\]

\begin{vidu}
Bếp A dùng LPG giá \(30\,000\) đồng/kg, nhiệt trị \SI{46}{\mega\joule\per\kilo\gram}, hiệu suất \(60\%\). Bếp B dùng một nhiên liệu khác giá \(18\,000\) đồng/kg, nhiệt trị \SI{30}{\mega\joule\per\kilo\gram}, hiệu suất \(75\%\). Hỏi nhiên liệu nào có chi phí thấp hơn để cung cấp \SI{100}{\mega\joule} nhiệt có ích?
\end{vidu}
\begin{loigiai}
Với LPG, năng lượng có ích từ \(1\) kg:
\[
46\times0.60=\SI{27.6}{\mega\joule}.
\]
Khối lượng cần để có \SI{100}{\mega\joule}:
\[
m_A=\frac{100}{27.6}\approx\SI{3.623}{\kilo\gram}.
\]
Chi phí:
\[
C_A\approx3.623\times30000\approx108700\text{ đồng}.
\]

Với nhiên liệu B:
\[
30\times0.75=\SI{22.5}{\mega\joule\per\kilo\gram},
\]
\[
m_B=\frac{100}{22.5}\approx\SI{4.444}{\kilo\gram},
\]
\[
C_B\approx4.444\times18000\approx80000\text{ đồng}.
\]

\textbf{Vậy,} để cung cấp cùng \SI{100}{\mega\joule} nhiệt có ích, nhiên liệu B có chi phí thấp hơn.
\end{loigiai}

\subsection{Năng lượng và lượng \texorpdfstring{\ce{CO2}}{CO2} phát thải}

Khi đề cho công thức nhiên liệu, có thể viết phương trình cháy để tính lượng \ce{CO2}. Nếu so sánh ``mức phát thải'' giữa hai nhiên liệu, cần xem đề yêu cầu so sánh theo cùng khối lượng, cùng thể tích hay cùng năng lượng có ích.

\begin{vidu}
Methane và propane cháy hoàn toàn:
\[
\ce{CH4 + 2O2 -> CO2 + 2H2O},
\]
\[
\ce{C3H8 + 5O2 -> 3CO2 + 4H2O}.
\]
Biết \(1\) mol methane tỏa \SI{890}{\kilo\joule}; \(1\) mol propane tỏa \SI{2220}{\kilo\joule}. So sánh số mol \ce{CO2} tạo ra khi mỗi nhiên liệu cung cấp \SI{4440}{\kilo\joule}, bỏ qua thất thoát.
\end{vidu}
\begin{loigiai}
Với methane:
\[
n_{\ce{CH4}}=\frac{4440}{890}\approx\SI{4.989}{\mole}.
\]
Mỗi mol methane tạo \(1\) mol \ce{CO2}, nên:
\[
n_{\ce{CO2}}^{(\ce{CH4})}\approx\SI{4.989}{\mole}.
\]

Với propane:
\[
n_{\ce{C3H8}}=\frac{4440}{2220}=\SI{2.00}{\mole}.
\]
Mỗi mol propane tạo \(3\) mol \ce{CO2}, nên:
\[
n_{\ce{CO2}}^{(\ce{C3H8})}=\SI{6.00}{\mole}.
\]

\textbf{Vậy,} để cung cấp cùng \SI{4440}{\kilo\joule}, methane tạo ít \ce{CO2} hơn propane theo các dữ kiện đã cho.
\end{loigiai}

\begin{tuluyen}
Một gia đình cần trung bình \SI{24}{\mega\joule} nhiệt có ích mỗi ngày. Bếp LPG có hiệu suất \(60\%\). LPG có nhiệt trị trung bình \SI{46}{\mega\joule\per\kilo\gram}. Tính khối lượng LPG dùng trong 30 ngày.
\dapso{\(\approx\SI{26.1}{\kilo\gram}\).}
\begin{loigiaiBT}
Tổng nhiệt có ích trong 30 ngày:
\[
Q_{\text{ích}}=24\times30=\SI{720}{\mega\joule}.
\]
Nhiệt lượng LPG phải tỏa:
\[
Q=\frac{720}{0.60}=\SI{1200}{\mega\joule}.
\]
Khối lượng LPG:
\[
m=\frac{1200}{46}\approx\SI{26.1}{\kilo\gram}.
\]
\textbf{Vậy,} gia đình cần khoảng \SI{26.1}{\kilo\gram} LPG trong 30 ngày.
\end{loigiaiBT}
\end{tuluyen}

\begin{tuluyen}
Một loại nhiên liệu có giá \(25\,000\) đồng/kg, nhiệt trị \SI{40}{\mega\joule\per\kilo\gram}; bếp sử dụng có hiệu suất \(65\%\). Tính số tiền trung bình phải trả để nhận được \SI{520}{\mega\joule} nhiệt có ích.
\dapso{\(500\,000\) đồng.}
\begin{loigiaiBT}
Mỗi kilogram nhiên liệu cung cấp nhiệt có ích:
\[
40\times0.65=\SI{26}{\mega\joule}.
\]
Khối lượng cần dùng:
\[
m=\frac{520}{26}=\SI{20}{\kilo\gram}.
\]
Chi phí:
\[
C=20\times25000=500000\text{ đồng}.
\]
\textbf{Vậy,} số tiền trung bình cần trả là \(500\,000\) đồng.
\end{loigiaiBT}
\end{tuluyen}
